% ARC237: LSPV-1 2026-10-29 19:00 America/Los_Angeles
\documentclass[10pt]{article}
\usepackage[margin=.8in]{geometry}
\usepackage{amsmath,amssymb,amsthm,xurl,hyperref}
\newtheorem{theorem}{Theorem}
\newtheorem{lemma}[theorem]{Lemma}
\title{ARC237: A Sharp Nonlinear-Preprocessing One-Probe Parity Frontier\\ Contextual Information Value and Construction--Storage--Continuation Costs}
\author{LSPV-1 Formal Proof Lead}
\date{29 October 2026, 19:00 PDT}
\begin{document}\sloppy\maketitle
\begin{center}\fbox{\parbox{.91\linewidth}{\textbf{PROVED/EXACT, RESTRICTED; independent referee pending.}
An arbitrary deterministic \emph{nonlinear} preprocessing map needs at least
$L=b-\lfloor\log_2(b+1)\rfloor$ persistent bits if all subsequent parity contexts
must be answered with at most one original-source sign-bit probe. A uniform
linear summary meets this bound for every $b$. This is a restricted
two-stage data-structure theorem, not a SAT runtime lower bound. P versus NP OPEN.}}\end{center}

\section{PROPOSE, SEARCH, and preceding work}
Target: Is the previously known project one-probe storage endpoint specific to
\emph{linear} preprocessing, or does it hold against all nonlinear summaries?
M1, a nonlinear summary beats the linear endpoint: FALSIFIED in this model.
M2, public syndrome cosets cover every future parity mask with one probe:
PROVED/EXACT.
M3, each one-bit observation has at most one unit of realized support-value gain:
FALSIFIED by a four-candidate asymmetric fiber.
M4, a certified polynomial-value rank and full per-stage charged costs yield
polynomial time: PROVED as an elementary conditional amortization criterion.

Relevant actual sources: repository AGENTS.md blob
\texttt{19656af102604d5f0a721d548cf2e4e72a063561}, CODEX.md blob
\texttt{754999f50de93cdb693ed37bc8110d31e9e698ef}, the research skill
blob \texttt{b8df83a7b0b1428f349b0e02ce1387d979e8db7f},
LSPV-2 CONTEXTUAL-VALUE-HAMMING235 blob
\texttt{b4a1a846dab97bd02691578d9abe5a081890d8b7}, its executable
blob \texttt{910d1a8380df58e2103ac54d2270cac0ed1ab5a2},
LSPV-4 HEAD-WINDOW-INFO-VALUE236 blob
\texttt{ace6000e9b8dbb52fe0a6487f4edf10a00648b89},
and root main.tex blob \texttt{4217fdab589157be0cccd4d021c591dd7356a8ac}
at preflight. Root already contains contextual interaction dividends
and detached-summary counting; these are not new ARC237 contributions.
External code inspected: sagemath/sage
\path{src/sage/coding/hamming_code.py} blob
\texttt{431e660fae77184147230223f66c76f2d5f04b38}
(Hamming parameters; Sage not executed), and QuMuLab/dsharp README
blob \texttt{a9a4e2e537c6b98377e07337e410c23218d05846}
(a CNF-to-d-DNNF compiler, not run). Darwiche--Marquis (2002)
distinguishes compilation size from query tractability; Hamming (1950)
is the coding-theory precedent. No claim of external mathematical priority.

\section{Exact source, quantifiers and charge model}
Fix $b\ge1$. The source $A_a(X)=\bigwedge_{i=1}^b(x_i=a_i)$
contains unknown source-sign vector $a\in\mathbb F_2^b$, in fixed-width
records of physical length $N=\Theta(b\log(b+2))$.
Preprocessing is \emph{any} deterministic total map
$E_b:\mathbb F_2^b\to\{0,1\}^L$, performed \emph{before}
the query mask $w\in\mathbb F_2^b$ and target parity bit $t$ are revealed.
The continuation is the explicit parity constraint $B_{w,t}=[w\cdot X=t]$
(with a charged XOR3-chain CNF encoding if desired);
$\mathrm{SAT}(A_a\wedge B_{w,t})\iff w\cdot a=t$.
All source-dependent persistent information, including head positions,
caches, tables and source-dependent advice, belongs in the exactly $L$
retained bits. The unchanged original source remains read-only,
accessible online through at most $Q$ named sign-bit probes.
The indexed bit-RAM charges those probes independently of arithmetic,
address formation, full offline construction $P$, peak work storage $S$
and physical input $N$. A fixed public decoder and public source layout
are permitted. No answer-class oracle or unchecked assignment is supplied.

\begin{theorem}[Nonlinear summary, one raw sign-bit probe]\label{th:one}
For every $b$, every source encoder $E_b$, and every total deterministic
decoder that correctly answers all $a,w,t$ with at most \emph{one}
post-query sign-bit probe, necessarily
\[
 \boxed{L\ge b-\lfloor\log_2(b+1)\rfloor}.
\]
A uniform deterministic linear encoder with exactly
$L=b-\lfloor\log_2(b+1)\rfloor$ bits achieves at most one probe
for all future parity masks, so the bound is sharp even among arbitrary
nonlinear encoders and nonlinear adaptive decoders.
\end{theorem}
\begin{proof}
For every nonempty summary fiber $F_s=E_b^{-1}(s)$, choose $a_0$ and
write $V_s=\operatorname{span}\{a+a_0:a\in F_s\}$, $r_s=\dim V_s$.
Then $|F_s|\le2^{r_s}$.
The parity functions $a\mapsto w\cdot a$ on $F_s$ form exactly
$2^{r_s}$ distinct classes \emph{modulo constants}: two masks agree up
to a constant iff their restrictions to $V_s$ agree, and every
functional on $V_s$ extends to $\mathbb F_2^b$.
Given $s$ and $w$, a one-probe deterministic decoder chooses whether
and where to read \emph{before} observing any source bit.
Its answer is therefore constant, $a_i$ or $1+a_i$.
Modulo constants these give at most $b+1$ distinct classes.
Exactness for every mask implies $2^{r_s}\le b+1$,
hence $|F_s|\le b+1$.
There are at most $2^L$ fibers, so
$2^b\le2^L(b+1)$ and the integer lower bound follows.

For achievability put $r=\lfloor\log_2(b+1)\rfloor$.
Choose a public rank-$r$ binary matrix $G$ whose $b$ columns contain
all $2^r-1$ nonzero $r$-bit vectors, with arbitrary duplicates in
remaining positions. By Gaussian elimination compute a basis
$k_1,\ldots,k_{b-r}$ for $\ker G$ and store only
$s=(k_1\cdot a,\ldots,k_{b-r}\cdot a)$, of length $L=b-r$.
For any future $w$, compute $Gw$. If it vanishes,
$w\in\ker G$ and $w\cdot a$ is recovered from $s$.
If not, find a column $Ge_i=Gw$, giving
$w+e_i\in\ker G$ and $w\cdot a=(w+e_i)\cdot a+a_i$.
The first term is a known combination of summary bits, and the
second uses one counted raw sign probe. The construction, elimination,
coefficient discovery, verification and output have polynomial bit cost.
\end{proof}

\section{Contextual value and complete cost accounting}
For a declared state $(s,K)$, finite candidate set $\mathcal H_{s,K}$,
target coordinates $W$ and score scale, define
\[
\Phi_W(s,K)=\log_2|\{\pi_W(a):a\in\mathcal H_{s,K}\}|,\qquad
v_\Phi(I\mid s,K)=\Phi_W(s,K)-\Phi_W(U_I(s,K)).
\]
Use $I\mid_{v_\Phi}$ and $(x=1)\mid_{v_\Phi}$ to denote a
contextually valued datum; the bar is neither probability conditioning
nor divisibility. Source acquisition, authentication, value discovery
and application costs $C_s(I)$ are distinct from this epistemic
support score. A redundant sound fact has value zero whenever the
scored state remains unchanged. Along any actual sequence, conditional
marginals telescope exactly; initial singleton labels need not add.

For target $t$ and auxiliary random bit $r$, facts
$I_1:t+r=1$ and $I_2:r=0$ each have zero initial target-support
value, but together determine $t$ and give unit synergy.
A repeated fact has zero marginal. Moreover, a possibility set
$\{000,100,010,001\}$ has $\log_2 4=2$ units of log-cardinality,
but observing one coordinate equal to $1$ can isolate one source
and drop this potential by two in \emph{one realized bit probe}.
Thus a universal per-execution one-bit progress cap is false.
Composite storage-penalized potentials may yield negative values,
without claiming negative Shannon information.

\begin{lemma}[Fully charged amortized upper-bound criterion]
Suppose a fixed uniform native-machine algorithm has full preprocessing
cost $P_0(N)$, and stages $j$ whose measured costs $c_j$ include
all failures, score computation, source access, joins, storage updates,
validation and output. Let an independently defined integer rank
$0\le R_0\le N^d$ decrease strictly at every stage, and let a
nonnegative, monotonically nonincreasing potential satisfy
$\Phi_0\le N^e$. If for every stage
$c_j\le A(N)(\Phi_j-\Phi_{j+1})+B(N)$,
where $A,B,P_0$ are polynomial, then
$T\le P_0+A\Phi_0+B R_0=\mathrm{poly}(N)$.
Peak space is polynomial only under a separate same-run workspace
invariant. Rank bit length alone, uncharged failures or assumed
SAT coverage do not meet the hypotheses.
\end{lemma}
\begin{proof}
There are at most $R_0$ stages. Sum their complete costs and
telescope the potential. No independent maxima from different
executions are combined.
\end{proof}

\section{Attainable restricted frontier, tests, remaining gap}
The optimal one-probe storage amount is
$L_\star=b-\lfloor\log_2(b+1)\rfloor$.
The uniform realization reads all $b$ source signs during offline
construction, runs Gaussian elimination in conservative $O(b^3)$
indexed bit operations with $O(b^2)$ temporary bits,
retains $L_\star$ bits, and answers a query with $Q\le1$
sign probe and $O(b^3)$ conservative bit operations.
The read-only original source still occupies $N$ bits; this
one-probe bound does not constrain single-tape head travel,
two-writable-tape writes or fixed-UTM native resources.
Controls: retaining $a$ uses $(L,Q)=(b,0)$;
retaining nothing and scanning requested signs uses worst-case
$(L,Q)=(0,b)$.
A joint solver that knows both $A,B$ before construction has
different quantifiers.

The independent executable
\path{research/agents/lspv-1/2026-10-29_1900_ARC237_nonlinear_one_probe.py}
tests every source and query mask for $b=1,\ldots,8$:
87,380 exhaustive pairs, all correct with no more than one
source probe. It independently enumerates all 256 one-bit
nonlinear preprocessors at $b=3$; exactly two work with one
probe, while no zero-bit summary works. For $b=2$,
six of sixteen one-bit maps succeed and the zero-bit map fails.
The raw JSON and full five-page local TeX/PDF have SHA manifests
in the corresponding reproducibility bundle.
Finite tests are COMPUTATIONAL, not theorem proofs.

\textbf{Next missing lemma:} For $Q=2$, bound the number
of parity restrictions representable on arbitrary nonlinear
fibers by adaptive depth-two source-bit decision trees, and
compare against the best linear covering-code construction.
Do not assume a universal SAT runtime gap.
Neither $P=NP$ nor $P\ne NP$ is proved.
\begin{thebibliography}{9}
\bibitem{DM} A. Darwiche and P. Marquis, A Knowledge Compilation Map,
JAIR 17 (2002), 229--264, DOI:10.1613/JAIR.989.
\bibitem{Ha} R. W. Hamming, Error Detecting and Error Correcting Codes,
Bell System Technical Journal 29 (1950), 147--160.
\bibitem{FS} M. Fredman and M. Saks, The Cell Probe Complexity of
Dynamic Data Structures, STOC (1989), 345--354.
\end{thebibliography}
\end{document}
